immini[thm]{\subsection*{I. Đặc điểm chuyển động của các phân tử khí}
\subsubsection{Chuyển động Brown trong chất khí}
\begin{itemize}
	\item Năm 1827, Robert Brown quan sát các hạt phấn hoa trong nước bằng kính hiển vi và nhận thấy chúng chuyển động không ngừng. 
	\item Ông ghi lại vị trí của hạt phấn hoa sau những khoảng thời gian xác định, rồi nối các điểm đó lại thì thu được một đường gấp khúc không theo trật tự nào.
	\item Khi nhiệt độ của nước tăng thì các hạt phấn hoa cũng chuyển động nhanh hơn.	
\end{itemize}}{includegraphics[width=4cm]{Images/brown}}
immini[thm]{\subsubsection{Giải thích của Einstein}
\begin{itemize}
	\item Các hạt phấn hoa bị các phân tử nước chuyển động hỗn loạn và đập vào từ mọi phía.
	\item Các phân tử nước có kích thước nhỏ hơn nhiều lần so với hạt phấn hoa và tại một thời điểm, số va đập trung bình của các phân tử nước vào hai phía đối diện bất kì của hạt phấn hoa không bằng nhau. Điều này gây nên chênh lệch áp lực và làm cho hạt phấn hoa chuyển động.
\end{itemize}
}
{
includegraphics[width=8cm]{Images/brown13T}
}
\begin{itemize}
	\item Hai phía đối diện có chênh lệnh áp lực là bất kì nên chuyển động của hạt phấn hoa không có phương ưu tiên.
	\item Nhiệt độ của nước càng cao thì các phân tử nước chuyển động càng nhanh và càng hỗn loạn hơn, nên hạt phấn hoa cũng sẽ chuyển động nhanh và hỗn loạn hơn.
\end{itemize}
\subsubsection{Chuyển động Brown của phân tử hạt khói lơ lửng trong không khí}
\Binhluan{
\begin{itemize}
	\item Ta cũng quan sát thấy chuyển động Brown của các hạt khói lơ lửng trong không khí. Hiện tượng này được giải thích tương tự như chuyển động của hạt phấn hoa trong nước. Trong đó: 
	\begin{itemize}
		\item Hạt khói đóng vai trò của các hạt phấn hoa.
		\item Phân tử không khí đóng vai trò của các phân tử nước.
	\end{itemize}
	\item Từ chuyển động Brown của các hạt khói trong không khí, có thể rút ra các kết luận sau:
	\begin{itemize}
		\item Chất khí được cấu tạo từ các phân tử chuyển động hỗn loạn, không ngừng.
		\item 	Nhiệt độ của khí càng cao thì tốc độ chuyển động hỗn loạn của các phân tử khí càng lớn.
\end{itemize}
\end{itemize}
}
{
Ở điều kiện chuẩn $[T=273\ K (0^\circ C)\ \text{và } p=1\ atm\ 10^5\ Pa)]$, các phân tử khí oxygen chuyển động với tốc độ trung bình vào khoảng 400 m/s.
}
\begin{note}
	Do các phân tử khí chuyển động hỗn loạn, liên tục va chạm với nhau và với thành bình nên tốc độ của chúng liên tục thay đổi. Vì vậy, nói đến tốc độ phân tử là nói đến \textit{tốc độ trung bình của các phân tử}\\
	\begin{align*}
		\overline{v}=\dfrac{v_1+v_2+...+v_n}{n}
	\end{align*}
\end{note}
\subsubsection{Tương tác giữa các phân tử khí}
Giữa các phân tử khí cũng có lực đẩy và lực hút, gọi chung là lực liên kết. Khoảng cách giữa các phân tử ở thể khí rất lớn so với ở thể lỏng và thể rắn nên lực liên kết giữa các phân tử ở thể khí rất yếu so với ở thể lỏng và thể rắn.
\subsection*{II. Mô hình động học phân tử chất khí}
\setcounter{paragraph}{0}
immini[thm]{\subsubsection{Mô hình}
\newghichu{
$\bullet$ Chất khí được cấu tạo từ các phân tử có kích thước rất nhỏ so với khoảng cách giữa chúng. Lực liên kết giữa các phân tử ở thể khí rất yếu so với ở thể lỏng và thể rắn.\\
$\bullet$ Các phân tử khí chuyển động hỗn loạn, không ngừng. Chuyển động này càng nhanh thì nhiệt độ của khí càng cao.\\
$\bullet$ Khi chuyển động hỗn loạn, các phân tử khí va chạm với nhau và với thành bình. Khi va chạm với thành bình thì các phân tử khí tác dụng lực, gây áp suất lên thành bình.\\

}}{includegraphics[width=4cm]{Images/boxt}}
\subsubsection{Khí lí tưởng}
\Binhluan{\newghichu{
$\bullet$ Các phân tử khí được coi là chất điểm, không tương tác với nhau khi chưa va chạm.
\begin{itemize}
	\item [$-$] Khi chưa va chạm. lực tương tác giữa các phân tử khí rất yếu, có thể bỏ qua.
	\item [$-$] Giữa hai va chạm liên tiếp, phân tử khí lí tưởng chuyển động thẳng đều.
\end{itemize}
$\bullet$ Các phân tử khí tương tác khi va chạm với nhau và với thành bình. Các va chạm này là va chạm hoàn toàn đàn hồi.
}
}{Khí lí tưởng là mô hình đơn giản của khí thực nhưng vẫn phản ánh được những đặc điểm cơ bản của khí thực}

